\scanpage{22}\mspage{7}
\uline{avec pour courbe universelle} (\uline{en coordonnées non homogò\`nos})
\[
\mathrm{Y}^{2}=4\,\mathrm{X}^{3}-\mathrm{g}_{2}\,\mathrm{X}-\mathrm{g}_{3}
\]
\margin{\handwritten{\uline{\unsure{section} unité} ?}}
\uline{et pour différentielle invariante} $\pi=\mathrm{d}\,\mathrm{X}/\mathrm{Y}$.

De plus, lorsque $2$ et $3$ sont inversibles, la donnée d'une forme différentielle invariante suffit à rigidifier une courbe de genre $1$.

Il résulte aussitôt de (2.2) qu'on a
\[
\mathrm{c}_{4}^{3}-\mathrm{c}_{6}^{2}=\underset{\inserted{\handwritten{12^{3}}}}{1728}\,\Delta
\tag*{(2.6)}
\]
et cette formule, se réduisant à une identité algébrique, est valable en toute caractéristique.

3.\qquad \uline{Cas où} $2$ \uline{et} $\mathrm{c}_{4}$ \uline{sont inversibles}.

Supposons que $\mathrm{E}$ soit une courbe de genre un sur une base $\mathrm{S}$ où $2$ soit inversible. La forme différentielle invariante $\pi$ étant donnée, on peut choisir $\mathrm{x}$ et $\mathrm{y}$ tels que dans (1.3) on ait $\mathrm{a}_{1}=\mathrm{a}_{3}=0$. Les changements de coordonnées respectant cette condition s'obtiennent en faisant $\mathrm{u}=1$, $\mathrm{s}=\mathrm{t}=0$ dans (1.2).

On a donc
\[
\mathrm{y}^{2}=\mathrm{x}^{3}+\mathrm{a}_{2}\,\mathrm{x}^{2}+\mathrm{a}_{4}\,\mathrm{x}+\mathrm{a}_{6}\,\raisebox{.14ex}{\scriptsize\unsure{,|'}}
\tag*{(3.1)}
\]
\[
\mathrm{b}_{2}=4\mathrm{a}_{2}\qquad\mathrm{b}_{4}=2\mathrm{a}_{4}\qquad
\mathrm{b}_{6}=4\mathrm{a}_{6}\qquad\mathrm{d}_{6}=8(\mathrm{a}_{2}\,\mathrm{a}_{4}-9\,\mathrm{a}_{6})
\]
\[
\mathrm{c}_{4}=2^{4}(\mathrm{a}_{2}^{2}-3\mathrm{a}_{4})\qquad
-\mathrm{c}_{6}=2^{6}\,\mathrm{a}_{2}^{3}-2^{5}\,3^{2}\,\mathrm{a}_{2}\,\mathrm{a}_{4}+2^{5}\,3^{3}\,\mathrm{a}_{6}\quad\text{et}
\tag*{(3.2)}
\]
\[
\Delta=2^{4}\,\delta\qquad,
\]
\clearpage
% archive sheet number: 21
